%%%%%%%%%%%%%%%%%%%%%%%%%%%%%%%%%%%%%%%%%
% Backtest of ICES Advice --- results briefing
% Official Imperial College London Beamer theme
% Affiliation: Sea++
%
% Compile with XeLaTeX from tex/ (Fonts/ and Images/ must resolve;
% they live in imperial/ --- see README note below):
%   xelatex beamer.tex
%
% The theme looks for Fonts/ and Images/ relative to the working
% directory. If compiling from tex/, either create junctions
%   Fonts -> imperial/Fonts , Images -> imperial/Images
% or compile from imperial/:
%   cd imperial && xelatex -output-directory=.. ../beamer.tex
%%%%%%%%%%%%%%%%%%%%%%%%%%%%%%%%%%%%%%%%%
%!TEX program = xelatex

\documentclass[
  aspectratio=169,
  t,
  onlytextwidth,
  10pt,
]{beamer}

% Find the Imperial theme when this file lives in tex/
\makeatletter
\def\input@path{{imperial/}}
\makeatother

\usetheme{Imperial}

\usepackage{amsmath,amssymb}
\usepackage{booktabs}
\usepackage{array}
\usepackage{colortbl}

\graphicspath{{imperial/Images/}{Images/}{figs/}{../figs/}}

\newcommand{\Fmsy}{F_{\mathrm{MSY}}}
\newcommand{\Ftar}{F_{\mathrm{target}}}
\newcommand{\Blim}{B_{\lim}}
\newcommand{\Btrig}{MSYB_{\mathrm{trigger}}}
\newcommand{\Bmsy}{B_{\mathrm{MSY}}}

\title{Backtest of ICES Advice}
\subtitle{Historical, $F_{\mathrm{target}}$, and Advice rule --- one SRR scenario}
\author{Laurence Kell}
\institute{Sea++}
\date{25 August 2026}

\begin{document}

%------------------------------------------------------------------------------
% Title
%------------------------------------------------------------------------------
\begingroup
  \setbeamercolor{background canvas}{bg=ICLBlue}
  \setbeamercolor{title page title}{fg=white}
  \setbeamercolor{title page subtitle}{fg=white}
  \setbeamercolor{author}{fg=white}
  \setbeamercolor{inst}{fg=white}
  \setbeamercolor{date}{fg=white}
  \setbeamertemplate{title page}[logo]{ICL_Logo_White.pdf}
  \frame[plain, s]{\titlepage}
\endgroup

%------------------------------------------------------------------------------
% Outline
%------------------------------------------------------------------------------
\begingroup
  \setbeamercolor{background canvas}{bg=ICLBlue}
  \setbeamercolor{normal text}{fg=white}\usebeamercolor[fg]{normal text}
  \setbeamercolor{page number in head/foot}{fg=white}
  \begin{frame}
    \begin{columns}[T]
      \begin{column}{0.42\linewidth}
        \HUGE\textbf{Outline}
      \end{column}
      \begin{column}{0.54\linewidth}
        \textbf{01}\tabto{0.125\linewidth} What was simulated\\
        \textbf{02}\tabto{0.125\linewidth} Methods\\
        \textbf{03}\tabto{0.125\linewidth} Two example stocks\\
        \textbf{04}\tabto{0.125\linewidth} Current (terminal year)\\
        \textbf{05}\tabto{0.125\linewidth} Future (20 years)\\
        \textbf{06}\tabto{0.125\linewidth} Caveats and next
      \end{column}
    \end{columns}
  \end{frame}
\endgroup

%------------------------------------------------------------------------------
% What was simulated
%------------------------------------------------------------------------------
\begin{frame}
  \frametitle{What was simulated}
  \framesubtitle{Six stocks, three trajectories, one Operating Model scenario}

  {\small
  \begin{tabular}{@{}llll@{}}
    \toprule
    Stock & Area & Assessment & Advice basis \\
    \midrule
    Irish Sea cod        & 7.a              & Age-structured & ICES MSY \\
    Irish Sea whiting    & 7.a              & SAM            & ICES MSY \\
    Celtic Sea cod       & 7.e--k           & SAM            & ICES MSY \\
    Celtic Sea whiting   & 7.b--c, 7.e--k   & Age-based      & ICES MSY \\
    Celtic Sea haddock   & 7.b--k           & SAM            & ICES MSY \\
    NE Atlantic mackerel & 1--8, 14, 9.a    & Integrated     & Coastal States MSY \\
    \bottomrule
  \end{tabular}}

  \vspace{0.45em}
  \begin{itemize}
    \item \textbf{Historical} --- latest WGCSE assessment.
    \item $\boldsymbol{F}_{\mathrm{target}}$ --- constant ICES $\Fmsy$ from 2005 (open-loop).
    \item \textbf{Advice rule} --- closed-loop \texttt{hcrICES}.
    \item \textbf{NE Atlantic mackerel} is a special case: a recent benchmark changed the assessment substantially, so the backtest mixes following advice with robustness to that change. Retained for completeness; removed from the final report.
  \end{itemize}

  Same Operating Model biology and SRR residuals; only the harvest policy differs.
  \textbf{Historical and Future = bh3}; ICES HCR parameters from \textbf{bh1}.
  Not the full Operating Model set.
\end{frame}

%------------------------------------------------------------------------------
% Methods
%------------------------------------------------------------------------------
\begin{frame}
  \frametitle{Methods}
  \framesubtitle{Open loop = simple projection; closed loop = MSE with feedback}

  \begin{columns}[T]
    \begin{column}{0.48\linewidth}
      \textbf{Open loop} --- $\boldsymbol{F}_{\mathrm{target}}$\\[0.25em]
      {\small Simple projection at constant $F_{\mathrm{target}}$.
      No feedback. \texttt{fwd} at 2005 SAG $\Fmsy$ from 2006 (lag 1).
      Benchmark: fish at $\Fmsy$ every year, without the HCR cutting $F$
      below $\Btrig$.}
    \end{column}
    \begin{column}{0.48\linewidth}
      \textbf{Closed loop} --- Advice rule\\[0.25em]
      {\small MSE with feedback control: the ICES HCR revises $F$ from
      SSB each year. \texttt{hcrICES}: $F=\Ftar$ when SSB$\ge\Btrig$;
      reduced toward $\Blim$. True Operating Model SSB (\texttt{err = NULL});
      \texttt{implErr = 0}; lag 1.}
    \end{column}
  \end{columns}

  \vspace{0.35em}
  {\small They can disagree --- Irish Sea cod historical
  $F \approx 0.07\times F_{\mathrm{target}}$.}

  \vspace{0.35em}
  {\small
  \begin{tabular}{@{}ll@{}}
    \toprule
    Setting & This draft \\
    \midrule
    Observation error & \texttt{err = NULL} (true Operating Model SSB) \\
    Implementation error & \texttt{implErr = 0} \\
    Advice lag & 1 year \\
    Stochasticity & Deterministic (iteration 1) \\
    HCR inputs & ICES SAG on bh1; held fixed \\
    Operating Model SRR & bh3 ($R_0$ prior; Historical and Future) \\
    \bottomrule
  \end{tabular}}
\end{frame}

\begin{frame}
  \frametitle{Control points $\neq$ performance metrics}
  \framesubtitle{Winker et al.\ 2026; Winker et al.\ 2025 (fsaf204)}

  \begin{itemize}
    \item ICES $\Ftar$, $\Blim$, $\Btrig$ are \textbf{operational control
          points}: they set $F$ in the HCR. Operating Model $\Bmsy$ and years-to-rebuild
          are \textbf{performance metrics} (CFP/MSY on the assumed truth).
    \item Hockey-stick derives those control points ($\Blim$ = breakpoint;
          $\Btrig \approx B_{\mathrm{pa}}$), not Operating Model truth. After the
          breakpoint recruitment is constant; the production function is
          truncated (Winker et al., \emph{Fish and Fisheries} 2026).
    \item $\Btrig$ is on average $\sim$60\% of $\Bmsy$. Using it as a
          surrogate over-classifies stocks as ``at MSY''.
    \item Why an Operating Model and closed loop: EqSim is largely an \textbf{open-loop}
          projection at constant $F$. It does not simulate MSE feedback;
          EqSim $\Fmsy$ is not precautionary relative to a shortcut MSE
          on an Operating Model (Winker et al., ICES \emph{JMS} 2025, fsaf204).
  \end{itemize}
\end{frame}

%------------------------------------------------------------------------------
% Conceptual figures
%------------------------------------------------------------------------------
\begin{frame}
  \frametitle{Two Historical stories}
  \framesubtitle{bh3 trajectories; ICES $\Btrig$ / $\Ftar$ from bh1 --- not 24 ggplot panels}

  \begin{columns}[T]
    \begin{column}{0.48\linewidth}
      \textbf{Irish Sea whiting --- rebuild}\\[0.35em]
      Historical $F/\Ftar \approx 5.3$.
      Constant $\Ftar$ (and AR once above trigger) cuts $F$ and rebuilds SSB to $\sim 5\times\Btrig$.

      \vspace{0.35em}
      {\small The figure you would look at: Historical SSB stays below trigger; $F_{\mathrm{target}}$ and AR rise well above it.}
    \end{column}
    \begin{column}{0.48\linewidth}
      \textbf{Irish Sea cod --- counter-example}\\[0.35em]
      Historical $F/\Ftar \approx 0.07$.
      Fishing at $\Ftar$ \textbf{lowers} terminal SSB versus history (AR $\Delta = -0.21$).

      \vspace{0.35em}
      {\small The figure you would look at: Historical $F$ already far below the ICES line; putting the rule in force is a \emph{rise} in $F$.}
    \end{column}
  \end{columns}

  \vspace{0.5em}
  Choke / recovery TAC versus an Operating Model / $\Ftar$ mismatch is not separated yet.
\end{frame}

%------------------------------------------------------------------------------
% Table-style findings
%------------------------------------------------------------------------------
\begin{frame}
  \frametitle{Current --- terminal year}
  \framesubtitle{SSB/$\Btrig$ and $F/\Ftar$ --- Historical bh3, HCR inputs ICES/bh1}

  {\footnotesize
  \begin{tabular}{@{}l rrrr r@{}}
    \toprule
    & Hist & $F_{\mathrm{target}}$ & AR & Hist & AR$-$Hist \\
    Stock & SSB/$\Btrig$ & SSB/$\Btrig$ & SSB/$\Btrig$ & $F/\Ftar$ & $\Delta$ SSB \\
    \midrule
    Irish Sea cod        & 0.54 & 0.27 & 0.33 & 0.07 & $-0.21$ \\
    Irish Sea whiting    & 0.51 & 4.94 & 4.95 & 5.25 & $+4.45$ \\
    Celtic Sea cod       & 0.02 & 4.22 & 4.23 & 4.44 & $+4.21$ \\
    Celtic Sea whiting   & 0.22 & 0.40 & 0.49 & 1.42 & $+0.27$ \\
    Celtic Sea haddock   & 0.63 & 1.21 & 1.21 & 1.78 & $+0.58$ \\
    NE Atlantic mackerel & 0.77 & 0.77 & 0.77 & 1.39 & $0.00$ \\
    \bottomrule
  \end{tabular}}

  \vspace{0.4em}
  All six below $\Btrig$ at the terminal year.
  Realised $F$ exceeds $\Ftar$ except Irish Sea cod.
  Under $F_{\mathrm{target}}$, $F/\Ftar = 1$ by construction.
\end{frame}

\begin{frame}
  \frametitle{Three findings to take away}
  \framesubtitle{Historical and Future (bh3) --- this scenario only}

  \begin{enumerate}
    \item \textbf{Irish Sea cod.} Fishing at $\Ftar$ lowers SSB versus history.
          Historical $F$ is already far below the ICES target.
    \item \textbf{Celtic Sea whiting.} ICES $\Btrig$ sits above Operating Model $\Bmsy$.
          A 20-year HCR can reach MSY biomass and still be ``below trigger''.
    \item \textbf{Celtic Sea cod.} From history, $\Btrig$ in $\sim$9 years;
          Operating Model $\Bmsy$ is not reached in 20 years from either start.
          The AR start is already above the trigger --- different targets.
  \end{enumerate}
\end{frame}

%------------------------------------------------------------------------------
% Future
%------------------------------------------------------------------------------
\begin{frame}
  \frametitle{Future --- 20-year HCR}
  \framesubtitle{bh3, mean biology, mean recruitment; both starts then follow the ICES HCR}

  {\footnotesize
  \begin{tabular}{@{}l rr rr@{}}
    \toprule
    & \multicolumn{2}{c}{From Historical} & \multicolumn{2}{c}{From Advice rule} \\
    \cmidrule(lr){2-3}\cmidrule(lr){4-5}
    Stock & $\Btrig$ & $\Bmsy$ & $\Btrig$ & $\Bmsy$ \\
    \midrule
    Irish Sea cod        & 8  & 6  & 8  & 7  \\
    Irish Sea whiting    & 2  & 7  & 0  & 0  \\
    Celtic Sea cod       & 9  & --- & 0  & --- \\
    Celtic Sea whiting   & --- & 18 & --- & 16 \\
    Celtic Sea haddock   & 4  & --- & 0  & --- \\
    NE Atlantic mackerel & 3  & 0  & 3  & 0  \\
    \bottomrule
  \end{tabular}}

  \vspace{0.4em}
  Years to threshold; $0$ = already above; --- = not within 20 years.
  Reaching ICES $\Btrig$ does not imply Operating Model $\Bmsy$.
  Mackerel is already above $\Bmsy$ while still below the (SAG) trigger.
\end{frame}

%------------------------------------------------------------------------------
% Caveats
%------------------------------------------------------------------------------
\begin{frame}
  \frametitle{Caveats}
  \framesubtitle{Read the numbers as one scenario, not a full MSE}

  \begin{itemize}
    \item \textbf{bh3.} Historical and Future trajectories use bh3
          ($R_0$ prior). HCR inputs (ICES $\Ftar$, $\Btrig$, $\Blim$) stay
          at SAG values on bh1. Operating Model shading ($F_{\mathrm{MSY}}$, $\Bmsy$)
          is bh3. bh1 is unconstrained Beverton--Holt --- steepness is
          estimated, not 1.
    \item \textbf{Mackerel $\Btrig$.} SAG field, equal to $B_{\mathrm{pa}}$ in
          the latest SAG row; not the LTMP / $\Blim$ advice-sheet convention.
    \item \textbf{No estimation error.} \texttt{err = NULL}: the HCR sees true Operating Model SSB.
          No implementation error; deterministic, not $P(\mathrm{SSB}{<}\Blim)$.
  \end{itemize}
\end{frame}

%------------------------------------------------------------------------------
% Next
%------------------------------------------------------------------------------
\begin{frame}
  \frametitle{Next}
  \framesubtitle{What this draft does not yet do}

  \begin{enumerate}
    \item \textbf{Other SRRs} --- bh2 (FishLife steepness) and Ricker rk1--rk3,
          with HCR inputs still held at ICES values.
    \item \textbf{Advised-catch series} --- to separate a choke / recovery TAC
          from an Operating Model / $\Ftar$ mismatch (Irish Sea cod).
    \item \textbf{Monte Carlo $P(\Blim)$} --- stochastic replicates for Historical
          and Future, not a single deterministic trajectory.
  \end{enumerate}
\end{frame}

%------------------------------------------------------------------------------
% Close
%------------------------------------------------------------------------------
\begingroup
  \setbeamercolor{background canvas}{bg=ICLBlue}
  \setbeamercolor{normal text}{fg=white}\usebeamercolor[fg]{normal text}
  \setbeamertemplate{closing slide logo}[logo]{ICL_Logo_White.pdf}
  \setbeamertemplate{closing slide text}[text]{Questions / discussion}
  \usebeamertemplate{closing slide}
\endgroup

\end{document}
